Origin-centred involution of \(c^{2}\) realised by four-quadrant paths of identical topology, labelled with the exact physical values of light-speed-squared

1. Foundational constants (exact)

\[
c = 299\,792\,458\,\mathrm{m\,s^{-1}} \quad\text{(exact by SI definition)}
\]

\[
c^{2} = c\cdot c = 89\,875\,517\,873\,681\,764\,\mathrm{m^{2}\,s^{-2}}
\]

\[
c^{2} = 8.9875517873681764\times 10^{16}\,\mathrm{m^{2}\,s^{-2}}
\]

\[
\frac{1}{c^{2}} = 1.1126500560536185\times 10^{-17}\,\mathrm{s^{2}\,m^{-2}}
\]

In the geometric construction these two numbers (and their signed counterparts) appear as the actual coordinate labels of the distinguished points. The visual height of \(c^{2}\) is normalised to unity so that the diagram remains legible; this normalisation is the self-referential scaling appropriate to units in which spacetime curvature (or the geometric role of \(c^{2}\)) is taken as the reference scale of 1.

2. The four paths – identical topology

Each path is a closed polygonal chain of four segments that begins and ends at the origin. The four paths are congruent under the action of the Klein four-group \(\{1,-1\}\times\{1,-1\}\) (independent sign flips of the two coordinates). Consequently they share exactly the same combinatorial and metric structure once absolute values are taken.

Top-right quadrant (prototype)  
\[
(0,0)
\;\longrightarrow\;
(\theta,\,c^{2})
\;\longrightarrow\;
\Bigl(\frac1\theta,\,c^{2}\Bigr)
\;\longrightarrow\;
\Bigl(\frac1\theta,\,\frac1{c^{2}}\Bigr)
\;\longrightarrow\;
(0,0)
\]

Top-left quadrant  
\[
(0,0)
\;\longrightarrow\;
(-\theta,\,c^{2})
\;\longrightarrow\;
\Bigl(-\frac1\theta,\,c^{2}\Bigr)
\;\longrightarrow\;
\Bigl(-\frac1\theta,\,\frac1{c^{2}}\Bigr)
\;\longrightarrow\;
(0,0)
\]

Bottom-right quadrant  
\[
(0,0)
\;\longrightarrow\;
(\theta,\,-c^{2})
\;\longrightarrow\;
\Bigl(\frac1\theta,\,-c^{2}\Bigr)
\;\longrightarrow\;
\Bigl(\frac1\theta,\,-\frac1{c^{2}}\Bigr)
\;\longrightarrow\;
(0,0)
\]

Bottom-left quadrant  
\[
(0,0)
\;\longrightarrow\;
(-\theta,\,-c^{2})
\;\longrightarrow\;
\Bigl(-\frac1\theta,\,-c^{2}\Bigr)
\;\longrightarrow\;
\Bigl(-\frac1\theta,\,-\frac1{c^{2}}\Bigr)
\;\longrightarrow\;
(0,0)
\]

(The parameter \(\theta>0\) is free; its concrete numerical value affects only the horizontal placement of the vertical segments, not the topological type.)

3. Geometric content of each segment

Outgoing radial ray  
   From the origin to \((\pm\theta,\pm c^{2})\).  
   This is the ordinary Euclidean ray that carries the full magnitude of \(c^{2}\).

Horizontal inversion segment  
   At constant height \(y=\pm c^{2}\), the abscissa is sent  
   \[
   \theta\;\longmapsto\;\frac1\theta.
   \]
   This is the restriction of the standard inversion \(x\mapsto 1/x\) to the horizontal line of height \(c^{2}\).

Vertical inversion segment  
   At constant abscissa \(x=\pm 1/\theta\), the ordinate is sent  
   \[
   c^{2}\;\longmapsto\;\frac1{c^{2}}.
   \]
   This is the restriction of the inversion \(y\mapsto 1/y\) to the vertical line of abscissa \(1/\theta\).

Returning radial ray  
   From \(\bigl(\pm 1/\theta,\pm 1/c^{2}\bigr)\) back to the origin.  
   The product of the two successive inversions (horizontal then vertical) is equivalent to the full plane inversion
   \[
   (x,y)\;\longmapsto\;\Bigl(\frac1x,\frac1y\Bigr)
   \]
   followed by a return to the origin along the inverted ray.

Because each path closes at the origin, the composition of the two inversions is realised as a closed loop based at the fixed point of the involution.

4. Why the construction is an involution

An involution is a transformation \(\iota\) satisfying \(\iota\circ\iota=\mathrm{id}\).  
The map
\[
\iota:(x,y)\;\mapsto\;\Bigl(\frac1x,\frac1y\Bigr)
\]
(defined on \(\mathbb{R}^{2}\setminus\{xy=0\}\)) is an involution: applying it twice recovers the original point.  

The polygonal paths realise this algebraic involution geometrically:
the outer points \((\pm\theta,\pm c^{2})\) are sent by \(\iota\) to the inner points \(\bigl(\pm 1/\theta,\pm 1/c^{2}\bigr)\);
the inner points are sent back by \(\iota\) to the outer points;
the origin is the unique fixed point (or, more precisely, the point at which the inversion is undefined yet is the centre of the whole construction).

The four sign combinations simply embed the same involution into every open quadrant, guaranteeing global consistency under the central inversion \((x,y)\mapsto(-x,-y)\).

5. Self-referential scaling and spacetime curvature

In geometric units (or Planck units, or any system in which \(c=1\)) one has \(c^{2}=1\).  
The diagram adopts precisely this convention for its vertical scale: the visual distance from the origin to the horizontal segments is declared to be 1.  
All numerical labels, however, remain the true SI values.  
This is the meaning of “self-referential scaling”: the curvature scale (or the conversion factor between mass and energy, or the coefficient that appears in the Einstein field equations when \(G\) is retained) is set to unity while the concrete laboratory numbers are retained as annotations.

6. Global symmetry group

The complete figure is invariant under:
reflection across the \(x\)-axis,
reflection across the \(y\)-axis,
rotation by \(180^\circ\) about the origin,
and therefore under the whole Klein four-group generated by these operations.

These symmetries are exactly the symmetries of an origin-centred involution that treats the two coordinates symmetrically after absolute values are taken.

7. Summary statement of the thesis

The four paths constitute a single geometric object: the origin-centred involution that interchanges the pair \(\{c^{2},1/c^{2}\}\) (and simultaneously the pair \(\{\theta,1/\theta\}\)) while preserving orientation of the axes up to sign.  
When the vertical scale is normalised so that the visual magnitude of \(c^{2}\) equals unity, the construction becomes the natural diagram of the self-referential geometric role played by light-speed-squared in spacetime geometry; the exact SI values appear only as labels on the vertices.  
The topology of every path is identical; only the quadrant (i.e., the choice of signs) distinguishes them.

The topology of every path is identical; only the quadrant (i.e., the choice of signs) distinguishes them.

Consider first the prototype path that occupies the open first quadrant. It begins at the origin, travels radially outward along the ray whose direction is fixed by a positive parameter \(\theta\) until it reaches the point whose ordinate equals the exact physical value of light-speed-squared, \(c^{2}=89\,875\,517\,873\,681\,764\,\mathrm{m^{2}\,s^{-2}}\). From that outer vertex a horizontal segment carries the abscissa through the inversion \(\theta\mapsto 1/\theta\) while the ordinate remains locked at \(+c^{2}\). A subsequent vertical segment then inverts the ordinate itself, sending \(c^{2}\) to its reciprocal \(1/c^{2}\approx 1.112650056\times 10^{-17}\,\mathrm{s^{2}\,m^{-2}}\), the abscissa now held fixed at \(+1/\theta\). The final leg returns radially to the origin, closing the loop. The resulting closed chain therefore consists of four edges and four vertices, and its combinatorial type is that of a simple quadrilateral pinched at the origin.

Now transport this same sequence of operations into the second quadrant by reversing the sign of every abscissa while leaving the ordinates untouched. The radial departure, the horizontal inversion, the vertical inversion, and the radial return occur in precisely the same order; only the signs of the \(x\)-coordinates have flipped. The number of edges remains four, the cyclic order of vertices is unchanged, and the local incidence relations at each vertex are identical. The path is therefore homeomorphic to the prototype via the continuous map \((x,y)\mapsto(-x,y)\).

The same reasoning applies to the lower half-plane. Reversing the sign of the ordinates alone maps the first-quadrant path onto the fourth-quadrant path; reversing both signs simultaneously maps it onto the third-quadrant path. In each case the transformation is a linear homeomorphism of the plane that preserves the combinatorial structure of the chain. Consequently every one of the four closed paths is topologically indistinguishable from the others: each is a simple closed polygonal loop based at the origin, each comprises exactly two radial edges and two inversion edges meeting at right angles, and each realises the same sequence of geometric operations—radial excursion, horizontal inversion, vertical inversion, radial return.

What alone distinguishes the four realisations is the open quadrant in which they reside, or, equivalently, the particular choice of signs attached to the pair \((\theta,c^{2})\). That choice of signs is not an alteration of topology; it is merely the embedding of one and the same abstract cell complex into each of the four connected components of \(\mathbb{R}^{2}\setminus(\{x=0\}\cup\{y=0\})\). The resulting global figure therefore possesses a discrete symmetry group isomorphic to \(\mathbb{Z}{2}\times\mathbb{Z}{2}\), generated by the two independent reflections that flip the signs. Under the action of this group the four paths are permuted among themselves while their intrinsic topological type remains invariant.

In this light the entire construction may be read as a single topological object—an origin-based involutive cell—replicated by sign changes across the coordinate axes. The physical magnitudes \(c^{2}\) and \(1/c^{2}\) appear only as labels on corresponding vertices; they do not affect the incidence structure. The parameter \(\theta\) likewise affects only the metric placement of the vertical edges, never their combinatorial role. Hence the statement that “the topology of every path is identical; only the quadrant distinguishes them” is not a loose analogy but a precise assertion of homeomorphic equivalence under the natural action of the sign group. This equivalence is what renders the four-quadrant diagram a coherent geometric embodiment of the origin-centred involution of light-speed-squared.
